\documentclass{beamer}
\usetheme{AnnArbor}
\usecolortheme{beaver}
\usefonttheme{default}
\usepackage[T1]{fontenc}
\usepackage[utf8]{inputenc}
\usepackage{lmodern}

\title{MegaUltraTweet}
\author{Balz Aschwanden}
\date{\today}

\begin{document}
\frame{\titlepage} 

\begin{frame}
  \frametitle{Logo} \pause
  \begin{center}
      \includegraphics[width=.4\textwidth]{megaultra_ruby}
  \end{center}
\end{frame}

\section{Team}
\begin{frame}
  \frametitle{Team} \pause
  \begin{minipage}{.2\textwidth}
    \includegraphics[height=.2\textheight]{raoul}
  \end{minipage}%
  \hfill
  \begin{minipage}{.8\textwidth}
    \textbf{Key Account Manager} \\ Raoul Grossenbacher \pause \\    
  \end{minipage}
  \begin{minipage}{.2\textwidth}
    \includegraphics[height=.2\textheight]{balz}
  \end{minipage}%
  \hfill
  \begin{minipage}{.8\textwidth}
    \textbf{Chief Deliverable Officer} \\ Balz Aschwanden \pause \\    
  \end{minipage}
  \begin{minipage}{.2\textwidth}
    \includegraphics[height=.2\textheight]{simon}
  \end{minipage}%
  \hfill
  \begin{minipage}{.8\textwidth}
    \textbf{Master Tracker} \\ Simon Curty \pause \\    
  \end{minipage}
  \begin{minipage}{.2\textwidth}
    \includegraphics[height=.15\textheight]{daniel}
  \end{minipage}%
  \hfill
  \begin{minipage}{.8\textwidth}
    \textbf{Quality Evangelist} \\ Daniel Ziltener \\    
  \end{minipage}
\end{frame}

\section{Problem}
\begin{frame}
  \frametitle{Problem} \pause
  Wie bleibe ich immer über die neusten Entwicklungen in einem Feld informiert? \pause
    \begin{itemize}
       \item Blogs \pause \textcolor{red}{Nicht immer relevant} \pause
       \item Nachrichtenmagazine \pause \textcolor{red}{Zu viele Neuigkeiten} \pause
       \item Netzwerk \pause \textcolor{red}{Auf persönlichen Kontakt angewiesen}
    \end{itemize}       
\end{frame}  

\begin{frame}
  \frametitle{Folge} \pause
  \begin{center}
    An wenige Quellen gebunden \pause\\
    % Ich will nicht täglich 100 Sachen in meinem RSS Feed
    Wirklich neue Entwicklungen gehen oft unter \pause\\
    % Wir bekommen sie nicht mit, weil wir zusehr damit beschäftigt sind, auf dem Laufenden zu bleiben
  \includegraphics[height=.5\textheight]{sadcrying}  
  \end{center}
\end{frame}
   
\section{Ziel}
\begin{frame}
  \frametitle{Ziel} \pause
  \begin{center}
    Grundlegend Neues sofort erfassen \pause \\
    Zugang zu bis jetzt unbekannten Quellen \pause \\
    \includegraphics[height=.4\textheight]{happy}     
  \end{center}
\end{frame}

\begin{frame}
  \frametitle{Gewinn für den User} \pause
  \begin{itemize}
    % Nur auf den relevanten von 50 Blogeinträgen
    \item Einfache Navigation im Begriffsfeld \pause
    \item Gezielte Auswahl nach Schlagwörtern \pause
    \item Direkten Zugriff auf Quellen \pause
    \item Sofortige Information bei wirklich neuen Entwicklungen
  \end{itemize}
\end{frame}

\section{Wissensbasis/Quelle}
\begin{frame}
  \frametitle{Wissensbasis/Quelle} \pause
  \begin{minipage}{.4\textwidth}
    \includegraphics[height=.4\textheight]{Twitter_logo_blue} \pause
  \end{minipage}%
  \hfill
  \begin{minipage}{.6\textwidth}
    \begin{itemize}
      \item Kurze Nachrichten \pause
      \item Sofortige Reaktion auf Neuerungen \pause
      \item Grosse Datenmenge
    \end{itemize}
  \end{minipage}
\end{frame}

\section{Technik}
\begin{frame}
  \frametitle{Technik} \pause
  \begin{minipage}{.3\textwidth}
    \includegraphics[height=.3\textheight]{RubyOnRails}
  \end{minipage}%
  \hfill
  \begin{minipage}{.7\textwidth}
    Ruby On Rails \pause \\    
  \end{minipage}
  \begin{minipage}{.3\textwidth}
    \includegraphics[height=.3\textheight]{MongoDB}
  \end{minipage}%
  \hfill
  \begin{minipage}{.7\textwidth}
    MongoDB \pause \\    
  \end{minipage}
  \begin{minipage}{.3\textwidth}
    \includegraphics[height=.3\textheight]{GitHub}
  \end{minipage}%
  \hfill
  \begin{minipage}{.7\textwidth}
    GitHub \\    
  \end{minipage}

\end{frame}

\section{Erste Iteration}
\begin{frame}
  \frametitle{Erste Iteration} \pause
  \begin{itemize}
    \item Twitter-Login für den User \pause
    \item Verbindung Thema / Hashtag \pause
    \item Schwerpunkt Thema \textit{Technik} \pause
    \item Mobile first
  \end{itemize}
  
\end{frame}

\begin{frame}
  \frametitle{Namensgebung} \pause
  \begin{center}
    Tweet bei Eintreten einer seltenen Kombination \pause \\
    \textbf{UltraTweet}
  \end{center}
\end{frame}

\begin{frame}
  \includegraphics[width=\textwidth]{BringItOnBro}
\end{frame}

\end{document}
